% ============================================================================
% Fichier  : partie4/P04_C11_transferts.tex
% Titre    : Les Transferts Internationaux
% Auteur   : MINKA MI NGUIDJOÏ Thierry Emmanuel
% Version  : 1.0 -- novembre 2026
% Corrigés : corriges/Corr_C11.tex
% ============================================================================
% Sources : loi n° 2024/017, art. 5 (transfert international, clauses
% contractuelles types), 21, 32, 53, 60, 69 (texte officiel) ; ISO/IEC 27018 ;
% recommandations 01/2020 du CEPD (comparaison) ; Minka Mi Nguidjoï,
% Massoussi Songo et Bialo (2026, manuscrit) pour la tension T3.
% Faits nouveaux (bible) : cartographie des flux de NOUN PAY, 21 au
% 25 septembre 2026.
% ============================================================================

\chapter{Les Transferts Internationaux}
\label{chap:transferts}
\chaptermark{Les transferts internationaux}

\epigraph{%
    « Une donnée hébergée à l'étranger a quitté le Cameroun,\\
    même si on la consulte depuis Bafoussam. »%
}{Principe du cours}

\niveauChapitre{Intermédiaire}{%
    Chapitres~\ref{chap:cartographie}, \ref{chap:loi} et~\ref{chap:soustraitants}
}

\begin{boxobjectifs}
\begin{itemize}
    \item Reconnaître un transfert international au sens de l'article~5,
          y compris l'hébergement et l'accès à distance.
    \item Appliquer les quatre vérifications de l'article~32.
    \item Cartographier les flux sortants d'une organisation à partir de
          sources indépendantes.
    \item Coter les transferts quand l'Autorité n'a publié ni liste de pays
          ni clauses types.
    \item Situer la tension entre transferts et coopération pénale.
\end{itemize}
\end{boxobjectifs}

\bigskip

% ============================================================================
\section*{Avant de commencer : « nos serveurs, c'est le nuage »}
% ============================================================================

Bafoussam, lundi 21~septembre~2026. Mme~KENFACK, responsable produit de NOUN
PAY, présente l'architecture de l'application à KORO AUDIT~: « Tout est dans
le nuage. » Où, exactement~? Elle cherche. L'hébergeur principal est
établi dans l'Union européenne~; la vérification d'identité par selfie est
confiée à un prestataire dont elle ne connaît pas le pays~; l'outil de
mesure d'audience et la messagerie d'entreprise sont américains~; le
support client de nuit est assuré par un centre d'appels au Maroc, qui
consulte les comptes à distance.

Aucun de ces flux n'a été autorisé. Mme~KENFACK pensait qu'un transfert
supposait d'« envoyer des fichiers à l'étranger »~; l'hébergement et la
consultation à distance ne lui semblaient pas en être.

\begin{boxalert}
\textbf{Violation typique.} Réduire le transfert à l'envoi de fichiers, et
oublier l'hébergement étranger, les services en ligne et la consultation
à distance depuis l'étranger.

\smallskip
\textbf{Règle~:} tout acheminement de données au-delà des frontières est un
transfert, quel qu'en soit le moyen. Il exige une autorisation préalable de
l'Autorité, après les vérifications de l'article~32.
\end{boxalert}

% ============================================================================
\section{Qu'est-ce qu'un Transfert}
\label{sec:c11-definition}
% ============================================================================

\begin{boxdef}
\termecle{Transfert international de données}~: selon l'article~5,
« acheminement des données à caractère personnel au-delà des frontières
d'un pays vers un autre, que ce soit vers un État tiers ou une organisation
internationale ».
\end{boxdef}

La définition vise l'\textit{acheminement}, pas le moyen. Sont des
transferts~: l'envoi d'un fichier à un partenaire étranger~; l'hébergement
d'une base chez un prestataire étranger~; l'utilisation d'un service en
ligne dont les serveurs sont à l'étranger (messagerie, formulaires, mesure
d'audience)~; la consultation à distance de données stockées au Cameroun
par une personne située à l'étranger, puisque les données traversent la
frontière pour s'afficher sur son écran. L'auditeur qui hésite sur un cas
le cote comme transfert et le signale~: la qualification restrictive
exposerait l'organisation aux articles~60 et~69.

% ============================================================================
\section{L'Article 32}
\label{sec:c11-article32}
% ============================================================================

\begin{boxjuris}
\textbf{\loicam{n°~2024/017}, article~32}~\cite{loi2024017}~:
\textit{« (1) Le transfert de données à caractère personnel vers un Etat
étranger, ou à une organisation internationale, est subordonné à
l'obtention d'une autorisation préalable délivrée par l'Autorité de
protection des données à caractère personnel dans les conditions
permettant de garantir l'exercice des droits de la personne concernée.
(2) A l'occasion de la délivrance de ladite autorisation, l'Autorité de
protection des données à caractère personnel doit préalablement
s'assurer~: que le pays de destination des données à caractère personnel
offre un niveau de protection suffisant~; de l'entrée en vigueur préalable
d'un instrument juridique signé avec le pays de destination des données à
caractère personnel transférées, en liaison avec les administrations et
organismes compétents~; de l'assujettissement de l'entité sollicitant
l'importation des données à caractère personnel aux règles contraignantes
de sécurité en matière de protection desdites données~; de la souscription
préalable, par les entités importatrices et exportatrices concernées, de
clauses contractuelles types en matière de transfert à l'international des
données à caractère personnel, édictées et publiées par l'Autorité de
protection des données à caractère personnel. (3) Les modalités de
délivrance de l'autorisation prévue à l'alinéa 1 ci-dessus sont fixées par
voie réglementaire. »}

\smallskip
\textbf{Portée.} Les quatre vérifications sont énumérées sans « ou »~: la
lecture littérale les rend cumulatives. Le deuxième point, un instrument
juridique « signé avec le pays de destination », est particulièrement
exigeant~; il suppose un accord d'État à État.
\end{boxjuris}

L'article~53 charge l'Autorité d'établir la liste des pays « reconnus comme
offrant un niveau de protection [\ldots] équivalent » et l'article~5
définit les \termecle{clauses contractuelles types} comme celles que
l'Autorité élabore et publie pour encadrer les transferts. Ni la liste ni
les clauses n'étaient publiées en novembre~2026.

\subsection{Coter sans liste ni clauses}

L'auditeur se trouve devant quatre conditions dont trois dépendent
d'instruments publics inexistants. Il décompose~:

\begin{itemize}
    \item « Le transfert est autorisé »~: \NON{} s'il n'existe pas
          d'autorisation, comme pour tout traitement.
    \item « Le pays de destination offre un niveau suffisant »~: \NEK{} de
          critère tant que la liste n'existe pas~; l'auditeur peut décrire,
          à titre d'information, le cadre du pays (une loi de protection,
          une autorité), sans en tirer de cote.
    \item « L'importateur est soumis à des règles contraignantes de
          sécurité »~: cotable sur pièces (contrat, certificat ISO/IEC~27018,
          rapport d'audit).
    \item « Les clauses contractuelles types sont souscrites »~: \NEK{} de
          critère tant qu'elles ne sont pas publiées.
    \item « L'organisation a documenté chaque transfert et préparé sa
          demande »~: cotable sur pièces.
\end{itemize}

\begin{boxinfo}
\textbf{Le modèle européen, en comparaison.} Le droit européen admet des
transferts fondés sur une décision d'adéquation, des clauses
contractuelles types adoptées par la Commission, des règles d'entreprise
contraignantes, ou des dérogations~; il demande d'évaluer, au cas par cas,
la législation du pays de destination~\cite{edpb2020transferts}. La loi
camerounaise ajoute à cet arsenal l'autorisation de l'Autorité et
l'instrument signé d'État à État. Les outils européens peuvent servir de
bonne pratique documentaire~; ils ne valent pas autorisation.
\end{boxinfo}

% ============================================================================
\section{Cartographier les Flux Sortants}
\label{sec:c11-flux}
% ============================================================================

On établit la liste des transferts à partir de trois sources
indépendantes~:

\begin{enumerate}
    \item \textbf{Les contrats et factures}~: chaque service payé à un
          prestataire étranger est un transfert probable~; le contrat dit
          parfois où sont les serveurs.
    \item \textbf{Les configurations et le code}~: adresses des services
          appelés par l'application, bibliothèques intégrées (mesure
          d'audience, vérification d'identité, notifications).
    \item \textbf{Les traces réseau}~: noms de domaine résolus et volumes
          échangés vers l'extérieur, rapprochés d'une base de localisation
          des adresses.
\end{enumerate}

La localisation d'une adresse IP est approximative~: un réseau de diffusion
de contenu répond depuis le point le plus proche, et un prestataire peut
stocker ailleurs qu'il ne répond. L'auditeur ne déduit pas le lieu de
stockage d'une adresse~; il le cherche dans le contrat ou la documentation
du prestataire, et cote \NEK{} quand il ne le trouve pas. Le laboratoire
APDC-LAB-07 exerce ce rapprochement.

% ============================================================================
\section{Les Flux de NOUN PAY}
\label{sec:c11-nounpay}
% ============================================================================

\begin{boxscenario}{Mission NOUN PAY. Six flux sortants, 21 au 25 septembre 2026}
\begin{tableaucours}{L{3.2cm}L{3.0cm}L{2.6cm}L{3.2cm}}{%
    \tableauHeader{Flux} &
    \tableauHeader{Données} &
    \tableauHeader{Pays} &
    \tableauHeader{Source de la preuve}
}
Hébergement principal & Toutes & Union européenne & Contrat, console \\
Vérification d'identité & Selfies, pièces d'identité & \NEK{} & Code de
l'application~; contrat muet \\
Mesure d'audience & Identifiants, parcours & États-Unis & Code, trafic \\
Messagerie d'entreprise & Échanges avec les clients & États-Unis &
Facture, documentation \\
Support de nuit & Consultation des comptes & Maroc & Contrat, journaux
d'accès \\
SMS & Numéros, messages & Cameroun & Contrat (pas un transfert) \\
\end{tableaucours}

\medskip
\textbf{Constatations.} Cinq transferts sur six flux sortants. Aucun n'est
autorisé, aucun n'est mentionné dans l'information des clients
(article~21). Le pays de traitement des selfies, données biométriques,
n'est établi par aucune pièce.

\textbf{Analyse} (confiance élevée). Le modèle d'entreprise de NOUN PAY
repose sur des services étrangers~: le rapatriement complet n'est pas
réaliste à court terme. La priorité est de documenter chaque transfert, de
préparer les demandes d'autorisation, d'informer les clients, et de lever
l'inconnu sur la biométrie.
\end{boxscenario}

\begin{boxlecture}{Les transferts de NOUN PAY}
\begin{tableaucours}{L{7.8cm}cL{3.2cm}}{%
    \tableauHeader{Exigence} &
    \tableauHeader{Cote} &
    \tableauHeader{Preuve}
}
Les transferts sont recensés & \OUI{} & Cartographie de l'audit \\
Les transferts sont autorisés & \NON{} & Aucune autorisation \\
Les clients sont informés des transferts & \NON{} & Notice de
l'application \\
Le pays de chaque transfert est connu & \NEK{} & Selfies, pays inconnu \\
Les importateurs sont soumis à des règles de sécurité & \OUI{} & Certificat
ISO/IEC 27018 de l'hébergeur \\
Le niveau de protection des pays est suffisant & \NEK{} & Liste non
publiée \\
Les clauses contractuelles types sont souscrites & \NEK{} & Clauses non
publiées \\
Les demandes d'autorisation sont préparées & \NON{} & Aucun dossier \\
\end{tableaucours}

\medskip
\textbf{Décompte.} Deux \OUI{}, trois \NON{}, trois \NEK{} sur huit.
\hfill\lecture{0{,}25}{0{,}375}{1{,}56}

\medskip
\textbf{Lecture.} $U = 0{,}375$~: le domaine est fortement opaque. Mais deux
des trois \NEK{} sont de critère et ne se lèveront qu'avec des textes~; un
seul est de preuve (le pays des selfies), et c'est lui qu'il faut lever en
priorité, parce qu'il porte sur des données biométriques.
\end{boxlecture}

% ============================================================================
\section{Transferts et Coopération Pénale}
\label{sec:c11-penal}
% ============================================================================

Les transferts ne concernent pas que les entreprises. Quand un officier de
police judiciaire camerounais demande à un fournisseur étranger de
conserver des données, ou quand une autorité étrangère demande au Cameroun
des données sur un suspect, des données personnelles franchissent une
frontière. La Convention de Budapest, à laquelle le Cameroun est Partie,
prévoit en son article~29 une procédure de conservation rapide entre
États. La loi n°~2024/017 soumet les transferts à autorisation préalable
(article~32), sous réserve des textes spécifiques à la justice
(article~72).

L'auteur de ce cours et ses coauteurs analysent cette rencontre comme une
« tension transfrontalière » dans un travail en cours~\cite{minka2026chainepenale}~:
elle n'est pas, selon eux, une incompatibilité, mais elle impose de
distinguer conservation, préservation, production, transmission et
transfert, qui peuvent relever de régimes différents. Le
Chapitre~\ref{chap:recherche} présente cette piste de recherche.

% ============================================================================
\section{Ce que les Transferts demandent à chaque Métier}
\label{sec:c11-metiers}
% ============================================================================

\begin{boxmetier}{Délégué à la protection des données}
Tenir une liste des transferts, avec pour chacun le pays, les données, le
prestataire, le contrat et l'état de la demande d'autorisation. Exiger de
chaque nouveau service en ligne la localisation de ses serveurs avant de
l'adopter.
\end{boxmetier}

\begin{boxmetier}{RSSI}
Fournir les traces des flux sortants et l'inventaire des bibliothèques
intégrées aux applications~; bloquer les services non déclarés.
\end{boxmetier}

\begin{boxmetier}{Juriste conformité}
Préparer les dossiers d'autorisation~: description du transfert, garanties
contractuelles, informations sur le pays, et, en attendant les clauses
types de l'Autorité, des clauses de protection équivalentes.
\end{boxmetier}

\begin{boxresume}
\begin{itemize}
    \item Tout acheminement de données au-delà des frontières est un
          transfert, y compris l'hébergement, les services en ligne et la
          consultation à distance.
    \item L'article~32 subordonne le transfert à une autorisation préalable,
          après quatre vérifications dont un instrument signé avec le pays
          de destination.
    \item Sans liste de pays ni clauses types publiées, deux de ces
          vérifications sont des \NEK{} de critère~; les autres se cotent sur
          pièces.
    \item Le lieu de stockage se prouve par le contrat ou la documentation,
          pas par la localisation d'une adresse IP.
    \item Le transfert sans autorisation est puni administrativement
          (article~60) et pénalement (article~69).
\end{itemize}
\end{boxresume}

\begin{boxinfo}
\textbf{Points de contrôle.} Vous avez compris ce chapitre si vous savez~:
\begin{itemize}
    \item dire si un service en ligne constitue un transfert~;
    \item énoncer les quatre vérifications de l'article~32~;
    \item distinguer un \NEK{} de critère d'un \NEK{} de preuve dans un audit
          des transferts.
\end{itemize}
\end{boxinfo}

\begin{boxexo}[title={\ding{45}\quad Exercice~11.1. Transfert ou non \niveaubase}]
(a) Un médecin de MBOLO consulte un dossier depuis un congrès à Paris~;
(b) NYANGA stocke ses sauvegardes chez un prestataire à Douala~; (c) NOUN
PAY envoie ses SMS par un opérateur camerounais~; (d) le logiciel bancaire
de NYANGA est hébergé à Abidjan~; (e) un formulaire de satisfaction en ligne
dont les serveurs sont aux États-Unis.
\end{boxexo}

\begin{boxexo}[title={\ding{45}\quad Exercice~11.2. Le dossier d'autorisation \niveaumoyen}]
Rédigez le dossier de demande d'autorisation du transfert d'hébergement de
NOUN PAY, rubrique par rubrique, en vous appuyant sur l'article~32.
Indiquez les pièces que vous joindrez.
\end{boxexo}

\begin{boxexo}[title={\ding{45}\quad Exercice~11.3. Les selfies \niveaumoyen}]
Proposez une démarche pour établir le pays de traitement des selfies.
Quelles pièces demandez-vous au prestataire~? Que recommandez-vous si
NOUN PAY ne parvient pas à l'établir~?
\end{boxexo}

\begin{boxexo}[title={\ding{45}\quad Exercice~11.4. Abidjan \niveauexpert}]
Le logiciel bancaire de NYANGA est hébergé en Côte d'Ivoire, qui dispose de
sa propre loi de protection des données. (a) Ce fait suffit-il à satisfaire
la première vérification de l'article~32~? (b) Quelles autres
vérifications restent à satisfaire~? (c) Comment cotez-vous la proposition
« le transfert vers Abidjan est conforme » au 30~septembre~2026~?
\end{boxexo}

\bigskip

\begin{boxinfo}
\textbf{Pour aller plus loin.}
\begin{itemize}
    \item Les recommandations 01/2020 du Comité européen de la protection des
          données~\cite{edpb2020transferts}, pour la méthode d'évaluation
          d'un pays de destination.
    \item L'ISO/IEC~27018~\cite{iso27018}, pour les obligations des
          hébergeurs.
\end{itemize}
\end{boxinfo}

\begin{boxinfo}
\textbf{Transition.} Une donnée qui sort sans autorisation est un écart~;
une donnée qui sort par effraction est une violation. Comment la détecter,
la qualifier et en informer qui de droit~? C'est l'objet du
Chapitre~\ref{chap:violations}.
\end{boxinfo}
